% GENERATED FILE. Edit canonical lessons, then run npm run generate:books.
% canonical-source-hash: fnv1a64:003193ce86c7a0a2
% canonical-lessons: HI-C223-adat, HI-C223-sauq, HI-C223-sapna, HI-C223-khelkud, HI-C223-ged

\chapter{A habit, A hobby, A dream, Sport, A ball}
\label{ch:hi-a-habit-a-hobby-a-dream-sport-a-ball}
\hlchaptermodality{223}

\begin{chapteropening}
\textbf{By the end of this chapter:} \emph{I can say a habit, a hobby, a dream, sport and a ball.}

Complete the last lesson of chapter 223: I can say a habit, a hobby, a dream, sport and a ball.
\end{chapteropening}

\section[\texorpdfstring{\hi{आदत} (ādat)}{आदत (ādat)}]{\hi{आदत} (ādat) — a habit}

\label{lesson:HI-C223-adat}



\begin{quote}
Before the new one: say the Hindi for a problem, then the Hindi for a remedy.
\end{quote}

\subsection*{You'll want to know: \hi{आदत}}
\textbf{\hi{आदत}} — \emph{ādat} — \textquotedblleft{}a habit\textquotedblright{}.

\begin{cousinweb}
\textbf{\hi{आदत}} (\emph{ādat}) is a loan from Arabic \emph{ʿāda}, \textquotedblleft{}custom, habit\textquotedblright{}, which reached Hindi through Persian. Its Arabic root is \emph{ʿ-w-d} (Lane).
\end{cousinweb}

\begin{grammarlens}[title={What you've built}]
One more word for this chapter.
\end{grammarlens}

\subsection*{Your turn}
\begin{itemize}
\raggedright
  \item \emph{Say it:} \emph{ādat}
  \item \emph{Say it:} \emph{ādat}, once more
  \item \emph{Say it:} say \emph{ādat}
  \item \emph{Recall:} say the Hindi for a problem, then the Hindi for a remedy, then say \emph{ādat} again
\end{itemize}

\begin{tcolorbox}[breakable,colback=teal!4,colframe=teal!35!black,title={Before you move on}]
What does \hi{आदत} mean? (\textquotedblleft{}A habit\textquotedblright{}.) Say it once more.
\end{tcolorbox}
\section[\texorpdfstring{\hi{शौक़} (śauq)}{शौक़ (śauq)}]{\hi{शौक़} (śauq) — a hobby, a liking}

\label{lesson:HI-C223-sauq}



\begin{quote}
Before the new one: say the Hindi for a remedy, then the Hindi for a habit.
\end{quote}

\subsection*{You'll want to know: \hi{शौक़}}
\textbf{\hi{शौक़}} — \emph{śauq} — \textquotedblleft{}a hobby, a liking\textquotedblright{}.

\begin{cousinweb}
\textbf{\hi{शौक़}} (\emph{śauq}) is a loan from Arabic \emph{shawq}, \textquotedblleft{}longing, desire\textquotedblright{}, which reached Hindi through Persian. Its Arabic root is \emph{sh-w-q} (Lane).
\end{cousinweb}

\begin{grammarlens}[title={What you've built}]
One more word for this chapter.
\end{grammarlens}

\subsection*{Your turn}
\begin{itemize}
\raggedright
  \item \emph{Say it:} \emph{śauq}
  \item \emph{Say it:} \emph{śauq}, once more
  \item \emph{Say it:} say \emph{śauq}
  \item \emph{Recall:} say the Hindi for a remedy, then the Hindi for a habit, then say \emph{śauq} again
\end{itemize}

\begin{tcolorbox}[breakable,colback=teal!4,colframe=teal!35!black,title={Before you move on}]
What does \hi{शौक़} mean? (\textquotedblleft{}A hobby, a liking\textquotedblright{}.) Say it once more.
\end{tcolorbox}
\section[\texorpdfstring{\hi{सपना} (sapnā)}{सपना (sapnā)}]{\hi{सपना} (sapnā) — a dream}

\label{lesson:HI-C223-sapna}



\begin{quote}
Before the new one: say the Hindi for a habit, then the Hindi for a hobby.
\end{quote}

\subsection*{You'll want to know: \hi{सपना}}
\textbf{\hi{सपना}} — \emph{sapnā} — \textquotedblleft{}a dream\textquotedblright{}.

\begin{cousinweb}
\textbf{\hi{सपना}} (\emph{sapnā}) goes back to Sanskrit \emph{svapna}, \textquotedblleft{}sleep, dream\textquotedblright{} (CDIAL 13904), though CDIAL does not list the Hindi form. The dream-words of the sister languages, such as Gujarati \emph{samṇũ}, sit under a separate form \emph{*supna}, \textquotedblleft{}dream\textquotedblright{} (CDIAL 13481), and Hindi is not among them either.
\end{cousinweb}

\begin{grammarlens}[title={What you've built}]
One more word for this chapter.
\end{grammarlens}

\subsection*{Your turn}
\begin{itemize}
\raggedright
  \item \emph{Say it:} \emph{sapnā}
  \item \emph{Say it:} \emph{sapnā}, once more
  \item \emph{Say it:} say \emph{sapnā}
  \item \emph{Recall:} say the Hindi for a habit, then the Hindi for a hobby, then say \emph{sapnā} again
\end{itemize}

\begin{tcolorbox}[breakable,colback=teal!4,colframe=teal!35!black,title={Before you move on}]
What does \hi{सपना} mean? (\textquotedblleft{}A dream\textquotedblright{}.) Say it once more.
\end{tcolorbox}
\section[\texorpdfstring{\hi{खेलकूद} (khelkūd)}{खेलकूद (khelkūd)}]{\hi{खेलकूद} (khelkūd) — sport, games}

\label{lesson:HI-C223-khelkud}



\begin{quote}
Before the new one: say the Hindi for a hobby, then the Hindi for a dream.
\end{quote}

\subsection*{You'll want to know: \hi{खेलकूद}}
\textbf{\hi{खेलकूद}} — \emph{khelkūd} — \textquotedblleft{}sport, games\textquotedblright{}.

\begin{cousinweb}
\textbf{\hi{खेलकूद}} (\emph{khelkūd}) is \textquotedblleft{}play-and-jump\textquotedblright{}: \emph{khel}, from a reconstructed \emph{*khel-}, \textquotedblleft{}play\textquotedblright{} (CDIAL 3918), and \emph{kūd}, from Sanskrit \emph{kūrdati}, \textquotedblleft{}leaps\textquotedblright{}, which CDIAL takes, probably, from Dravidian (CDIAL 3412).
\end{cousinweb}

\begin{grammarlens}[title={What you've built}]
One more word for this chapter.
\end{grammarlens}

\subsection*{Your turn}
\begin{itemize}
\raggedright
  \item \emph{Say it:} \emph{khelkūd}
  \item \emph{Say it:} \emph{khelkūd}, once more
  \item \emph{Say it:} say \emph{khelkūd}
  \item \emph{Recall:} say the Hindi for a hobby, then the Hindi for a dream, then say \emph{khelkūd} again
\end{itemize}

\begin{tcolorbox}[breakable,colback=teal!4,colframe=teal!35!black,title={Before you move on}]
What does \hi{खेलकूद} mean? (\textquotedblleft{}Sport, games\textquotedblright{}.) Say it once more.
\end{tcolorbox}
\section[\texorpdfstring{\hi{गेंद} (gẽd)}{गेंद (gẽd)}]{\hi{गेंद} (gẽd) — a ball}

\label{lesson:HI-C223-ged}



\begin{quote}
Before the new one: say the Hindi for a dream, then the Hindi for sport.
\end{quote}

\subsection*{You'll want to know: \hi{गेंद}}
\textbf{\hi{गेंद}} — \emph{gẽd} — \textquotedblleft{}a ball\textquotedblright{}.

\begin{cousinweb}
\textbf{\hi{गेंद}} (\emph{gẽd}) is a tadbhava, a word inherited from Sanskrit and reshaped by sound change. It comes from \emph{genduka}, \textquotedblleft{}ball for play\textquotedblright{} (CDIAL 4248). CDIAL notes that the related \emph{kanduka} has been traced to Dravidian.
\end{cousinweb}

\begin{grammarlens}[title={What you've built}]
One more word for this chapter.
\end{grammarlens}

\subsection*{Your turn}
\begin{itemize}
\raggedright
  \item \emph{Say it:} \emph{gẽd}
  \item \emph{Say it:} \emph{gẽd}, once more
  \item \emph{Say it:} say \emph{gẽd}
  \item \emph{Recall:} say the Hindi for a dream, then the Hindi for sport, then say \emph{gẽd} again
\end{itemize}

\begin{tcolorbox}[breakable,colback=teal!4,colframe=teal!35!black,title={Before you move on}]
What does \hi{गेंद} mean? (\textquotedblleft{}A ball\textquotedblright{}.) Say it once more.
\end{tcolorbox}
